\chapter{Case Study: Building a Linker with Lua Memory Maps}
\label{ch:lua-linker}

The assembler case study ends with relocatable objects whose addresses remain undecided.
The linker in \path{example-code/lua-toolchain/luald.c} completes the path: it reads those objects, resolves definitions, applies Lua-directed placement, patches relocations, and emits an executable ELF image.
The final test is a native x86-64 Linux process, but the verification also inspects the intermediate addresses and fields.
An exit status alone would not establish that the image's unused pointer or zero-fill metadata is correct.

The implementation embeds Lua 5.4 in C11 and emits ELF bytes directly.
It uses no external linker to construct its image.
External tools provide independent producers and consumers in the tests, while a separate host-loader experiment covers dynamic linking outside this static implementation.
The supported scope is an executable contract rather than an implied implementation of every feature discussed in the four linking chapters.

\section{Define the input and output profiles}
\label{sec:ch60-profile}

Input objects must be little-endian x86-64 ELF64 relocatable files with an ordinary section table and one static symbol table.
Allocated input sections may contain ordinary bytes or NOBITS storage and may request allocation, writing, or execution flags.
The linker accepts local/global/weak symbols with the supported ordinary types, default or hidden visibility, ordinary defining sections, and bounded absolute definitions.
It supports RELA relocations of types 64, PC32, PLT32, unsigned 32, and signed 32.

Archives, section groups, common symbols, TLS sections, dynamic objects, and unsupported relocation kinds are rejected.
Nonallocated debug-style sections are not copied into the output, and relocations targeting such sections are not applied.
The output is a static \code{ET\_EXEC} image containing program headers and loadable payloads, without a section-header table or dynamic-linking metadata.
It supports the tested syscall-only startup programs rather than a general hosted C runtime.

Table~\ref{tab:ch60-capacities} records implementation capacities.
The total output buffer is bounded as well as individual region capacities, so multiple individually valid region requests can still exceed the output limit.
That outcome is diagnosed rather than allowing unchecked allocation growth.

\begin{table}[htbp]
\caption{Explicit capacities bound the teaching implementation without silently changing accepted object semantics.}
\label{tab:ch60-capacities}
\centering
\begin{tabular}{ll}
\toprule
Resource & Limit \\
\midrule
Input objects & $64$ \\
Sections per object & $64$ \\
Symbols per object & $4096$ \\
Relocations per RELA section & $8192$ \\
Lua memory regions & $16$ \\
Placement rules & $64$ \\
Individual file/buffer or region extent & $16$ MiB \\
Supported alignment & power of two, at most $4096$ \\
\bottomrule
\end{tabular}
\end{table}

\section{Treat the object reader as a typed graph validator}
\label{sec:ch60-reader}

The ELF reader first checks the identification bytes, class, byte order, machine, version, file type, and record sizes.
Only then does it interpret section counts and offsets.
It verifies the section-header range before reading individual headers and requires a zero null header under this nonextended-count profile.
For every file-backed section it checks that offset is within the file and size is no greater than the remaining bytes.
This subtraction-based check avoids overflow in an unchecked \code{offset+size} expression.

The section-name string table must itself be a valid string table.
Every name offset must point inside its payload, and a terminating zero must occur before the end.
The same rule applies independently to symbol names in the symbol string table.
Checking only the name offset is insufficient: an unterminated string could otherwise make a later comparison read into unrelated records.

The symbol reader checks entry size, divisibility, capacity, local-before-global ordering, binding/type/visibility support, and defining-section bounds.
An ordinary symbol's value and size must fit within its defining section.
Absolute and undefined entries use distinct paths.
The null symbol is required to have the specified zero fields.
Only validated indices are retained for later resolution.

\begin{algorithm}[htbp]
\caption{Validate an object before exposing its references to layout or relocation.}
\label{alg:ch60-reader}
\begin{pseudocode}
\Require file byte buffer with known length
\Ensure validated sections, symbols, and relocation-table links
\State check ELF identification, target profile, header widths, and table bounds
\ForAll{section headers}
  \State validate size, alignment, file extent where present, and supported allocated semantics
\EndFor
\State validate section-name strings through the designated string table
\State locate one symbol table and validate its entry and string-table contracts
\ForAll{symbol entries}
  \State validate binding, type, visibility, index, and defining extent
\EndFor
\ForAll{RELA sections}
  \State validate symbol-table link, target-section index, entry width, and count
\EndFor
\end{pseudocode}
\end{algorithm}

Algorithm~\ref{alg:ch60-reader} establishes structural facts used by later passes.
Field-width and relocation-symbol checks still occur when each relocation is applied, because its kind determines how many bytes are required.
Figure~\ref{fig:ch60-validation} follows the layers of validation.
An error at one layer prevents the next from interpreting unvalidated numeric fields as trusted identities.

\begin{figure}[htbp]
\centering
\begin{tikzpicture}[node distance=13mm]
\node[basic block] (bytes) {file bytes and actual length};
\node[basic block,below=of bytes] (headers) {checked headers and extents};
\node[basic block,below=of headers] (names) {checked strings and symbols};
\node[basic block,below=of names] (rels) {checked relocation links and fields};
\draw[cfg edge] (bytes)--(headers); \draw[cfg edge] (headers)--(names); \draw[cfg edge] (names)--(rels);
\end{tikzpicture}
\caption{The reader validates the object's reference graph before using it to construct a memory image.}
\label{fig:ch60-validation}
\end{figure}

\section{Execute the map, then copy its policy into C}
\label{sec:ch60-map}

The map script returns one table with \code{entry}, \code{regions}, and \code{sections} fields.
Listing~\ref{lst:ch60-map} is the complete map used by the demo.
Lua arithmetic computes capacities from a page constant, while dense arrays determine declaration and placement order.
No hash-table iteration order is used as an implicit layout policy.

\begin{listing}[htbp]
\inputminted{lua}{example-code/lua-toolchain/map.lua}
\caption{The Lua map supplies explicit region boundaries and an ordered exact-name placement policy.}
\label{lst:ch60-map}
\end{listing}

The host requires plain tables without metatables and dense integer-indexed arrays.
This keeps configuration access from invoking hidden Lua callbacks after the protected script call has completed.
String fields are copied with checked lengths and no embedded NUL.
Address, length, and alignment fields must be nonnegative Lua integers.
The region bases and lengths are page-aligned, lengths are positive and bounded, and their endpoints remain inside the supported low-half address range.

Region names must be unique, region extents must not overlap, and each section rule must name an existing region.
Exact section-name rules must also be unique.
The implementation accepts region permissions \code{r}, \code{rx}, or \code{rw}; an input section requiring writing or execution must be placed in a region that permits that access.
Lua selects policy, but C performs every placement and range check.
The Lua manual defines the integer and table behavior used at this boundary~\cite{ierusalimschy-2026-lua}.

\section{Placement order and segment order are different orders}
\label{sec:ch60-order}

Placement visits section rules in their array order.
Within one rule it visits input objects in command-line order and their matching sections in section-table order.
This defines a reproducible concatenation policy.
It records, for every mapped input section, its region, offset within that region, final virtual address, and whether the section contributes file bytes or only memory extent.

Program-header emission has an additional ELF requirement: loadable segments appear in increasing virtual-address order~\cite{xinuos-2026-elf}.
The linker therefore sorts a small array of region indices for emission.
It does not reorder the region records themselves after section placements have stored their region identifiers.
Reordering the backing array would invalidate those identifiers unless every reference were translated too.
Figure~\ref{fig:ch60-order} separates the two orders.

\begin{figure}[htbp]
\centering
\begin{tikzpicture}[node distance=17mm and 22mm]
\node[basic block] (rules) {section-rule order};
\node[basic block,right=of rules] (objects) {input-object order};
\node[basic block,below=of rules] (place) {input placement records};
\node[basic block,below=of objects] (sort) {region index permutation};
\node[basic block,below=of sort] (ph) {ascending PT\_LOAD addresses};
\draw[cfg edge] (rules)--(place); \draw[cfg edge] (objects)--(place);
\draw[cfg edge] (sort)--(ph); \draw[dashed,->] (place)--node[below=3mm,font=\small] {stable region IDs}(sort);
\end{tikzpicture}
\caption{Placement preserves declared contribution order while segment emission uses a separate address-sorted permutation.}
\label{fig:ch60-order}
\end{figure}

The test suite deliberately declares the lower-address region after the higher-address region.
It verifies sorted loadable headers and executes the resulting image.
This test catches the tempting but incorrect implementation that simply emits program headers in Lua declaration order.

\section{Resolve symbols through their owning objects}
\label{sec:ch60-symbols}

A local relocation target resolves directly through the referencing object's symbol entry and placed section.
A nonlocal target invokes the strong/weak selection rule from Chapter~\ref{ch:link-symbols}.
A selected ordinary definition becomes its placed section address plus symbol offset; an absolute definition uses its stored value.
An unresolved weak name becomes zero, and an unresolved strong name produces a diagnostic.

The implementation scans for definitions rather than constructing an optimized global index.
If $T$ is the total number of symbols and $R$ the number of applied relocations, repeated full scans give a coarse $O(T^2+RT)$ resolution bound.
Section matching similarly costs $O(QH)$ for $Q$ placement rules and $H$ input sections, before copying payloads.
These bounds explain the deliberate capacities and identify the first scalability improvements: indexed exact-name placement and cached global resolution.

Tests link two objects that both define a private label with the same spelling.
They also compare weak-only resolution with weak-plus-strong resolution in both object orders.
The resulting program returns the value from the selected definition, and a repeated strong object produces a duplicate-definition error.
This exercises symbol policy through executable behavior as well as table inspection.

\section{Apply relocations without hidden truncation}
\label{sec:ch60-relocations}

Each allocated relocation target has a byte-mark array recording fields already patched.
The array belongs to the target input section, not to one relocation section, so two RELA sections cannot evade the overlap check by targeting the same field independently.
Before touching output bytes, the linker validates the relocation type, field extent, symbol index, and arithmetic range.
It refuses relocations into NOBITS storage because there is no initialized payload field to patch under this profile.

The PC32/PLT32 path subtracts the place address before checked addition of the explicit addend.
The unsigned-32 path requires a result between zero and \code{UINT32\_MAX}.
The signed-32 path requires \code{INT32\_MIN} through \code{INT32\_MAX}.
The full-width 64-bit path uses the specified unsigned bit arithmetic.
All writes use the explicit little-endian byte helper.

\begin{algorithm}[htbp]
\caption{Link the complete static image in dependency order.}
\label{alg:ch60-pipeline}
\begin{pseudocode}
\Require Lua map and ordered validated-profile relocatable inputs
\Ensure an executable image and deterministic placement report, or an error
\State execute Lua once and validate/copy its returned map
\State read each object and validate its structural references
\State assign every nonempty allocated section to an explicit placement rule
\State validate capacities, permissions, alignment, and initialized-prefix ordering
\State resolve global definitions and all required undefined references
\State apply checked relocations using final placements
\State verify the entry names a byte inside an executable contribution
\State sort loadable region indices by virtual address and emit program headers
\State write the completed image, set executable mode, and report placements
\end{pseudocode}
\end{algorithm}

Algorithm~\ref{alg:ch60-pipeline} makes final layout precede relocation.
The current implementation performs no later transformation that changes section sizes or addresses.
If relaxation or dead-section removal were added, they would need to run before this final patch phase or explicitly invalidate its results.
Keeping this ordering visible is more reliable than scattering address adjustments across instruction-specific patch code.

\section{Emit an initialized prefix and a zero-filled tail}
\label{sec:ch60-image}

Each region owns a buffer for its initialized prefix and a separate memory extent.
Initialized contributions are copied at their placed offsets, with intervening initialized gaps zero-filled.
Once a nonempty NOBITS contribution is encountered, later nonempty initialized contributions in that region are rejected.
This ensures the region can be represented by one file-backed prefix followed by a loader-supplied zero tail.

The writer aligns each region payload's file offset to 4096 and emits a \code{PT\_LOAD} header containing its virtual base, file size, memory size, flags, and alignment.
The entry is the selected \code{\_start} definition and must lie strictly inside an executable input section rather than merely equal its end.
The image also carries a non-executable stack descriptor.
The writer checks all semantic conditions before opening the output path; detected write failures remove the partial file.

Figure~\ref{fig:ch60-pipeline} shows the two independent end-to-end paths used for the demo.
Our assembler's object is consumed by both our linker and GNU \code{ld}.
Our linker also consumes a restricted object produced by GNU \code{as}.
These paths reduce the chance that a shared private misunderstanding between our writer and reader masks an object-format error.

\begin{figure}[htbp]
\centering
\begin{tikzpicture}[node distance=16mm and 22mm]
\node[basic block] (as) {luaas};
\node[basic block,right=of as] (obj) {ELF64 objects};
\node[basic block,below left=of obj] (own) {luald + Lua map};
\node[basic block,below right=of obj] (gnu) {GNU ld};
\node[basic block,below=of own] (a) {native exit 42};
\node[basic block,below=of gnu] (b) {native exit 42};
\draw[cfg edge] (as)--(obj); \draw[cfg edge] (obj)--(own); \draw[cfg edge] (obj)--(gnu);
\draw[cfg edge] (own)--(a); \draw[cfg edge] (gnu)--(b);
\end{tikzpicture}
\caption{Independent consumers check the teaching assembler's object contract through two native execution paths.}
\label{fig:ch60-pipeline}
\end{figure}

\section{Run the image and inspect its structure}
\label{sec:ch60-run}

After the assembler commands in Chapter~\ref{ch:lua-assembler}, invoke the linker:
\begin{verbatim}
/tmp/opencode/lua-tools/luald example-code/lua-toolchain/map.lua \
  /tmp/opencode/lua-tools/demo /tmp/opencode/lua-tools/demo.o \
  /tmp/opencode/lua-tools/answer.o
readelf -h -l /tmp/opencode/lua-tools/demo
/tmp/opencode/lua-tools/demo
\end{verbatim}
The final command exits with status $42$, intentionally nonzero as a program result.
The map report names each nonempty input contribution, its address, and size.
It is derived from the same validated placements used for relocation, so it explains the actual image rather than a separately recomputed approximation.

The hello example instead writes a Lua-generated message and exits with status zero.
Its RIP-relative address-forming instruction continues to work when the code region is rebased, because the message and referencing instruction move together.
The tests link and execute it at two different code origins.
This verifies a specific position-relative layout change, not a claim that the output executable is an arbitrary-base PIE.

\section{Test both success and rejection boundaries}
\label{sec:ch60-tests}

Table~\ref{tab:ch60-tests} organizes the end-to-end checks by the contract they test.
The tool builds run under GCC and Clang, each normally and with address/undefined-behavior sanitizers and leak detection.
The separate dynamic-loader/TLS experiment is compiled and executed with both compilers using the host loader.
All generated objects, executables, maps, and malformed fixtures live in a temporary directory outside the repository.

\begin{table}[htbp]
\caption{The verification combines independent arithmetic, production-tool interoperability, and explicit failure cases.}
\label{tab:ch60-tests}
\centering
\begin{tabular}{p{0.30\linewidth}p{0.57\linewidth}}
\toprule
Contract & Executed check \\
\midrule
Instruction encoding & Compare 276 selected forms against GNU assembler bytes. \\
Object interoperability & GNU linker consumes luaas objects; luald consumes restricted GNU objects. \\
Relocation arithmetic & Independently check both PC32 fields and the absolute demo pointer. \\
Memory image & Inspect program headers and execute a GNU-produced BSS-reading program. \\
Symbol policy & Local-name separation, weak/strong order, missing and duplicate definitions. \\
Layout policy & Reordered regions, rebased code, overlap, alignment, permission, and missing-rule rejection. \\
Range and parsing & Exact PC32 endpoints, overflow, truncated ELF, bad indices, and unsupported kinds. \\
\bottomrule
\end{tabular}
\end{table}

Recent GNU assemblers can emit allocated processor-property notes by default.
Those notes are outside the teaching linker's accepted profile.
The interoperability fixture explicitly assembles with \code{-mx86-used-note=no}; the linker does not silently discard unknown allocated metadata.
This makes the producer configuration part of the test contract instead of hiding a compatibility gap.

For PC32 boundaries, the test constructs a valid object and adjusts its explicit addend so that the mathematical result is exactly $-2^{31}$, $2^{31}-1$, or one unit outside either endpoint.
The two endpoints must serialize correctly, and the outside values must fail.
These images are inspected rather than executed, because a range-valid branch field alone does not establish that the deliberately chosen target is mapped code.
Separating format acceptance from safe execution of a particular fixture is essential to a meaningful test.

\section{Extensions need their own invariants}
\label{sec:ch60-extensions}

Archive support would add demand-driven object selection before the current resolution/layout stages.
Section garbage collection would add a sound root/edge graph and discard decisions before final placement.
Thunk insertion would add a distance-sensitive layout iteration.
Dynamic linking would add runtime tables and a loader protocol rather than a few extra relocation cases.
Each extension should preserve the existing narrow contract until its new behavior is explicitly implemented and tested.

A production tool would also replace global process-lifetime state with per-link contexts, indexed lookup, richer source diagnostics, atomic file replacement, and more complete metadata preservation.
Those engineering changes matter for scale and embedding.
The present artifact establishes the smaller but complete path from Lua-generated assembly through interoperable objects to a Lua-mapped executable.
The accompanying theory explains how larger tools extend that path without suggesting those extensions are already present in the case study.

\section{The emitted image as input to another compiler}
\label{sec:ch60-lifting-link}

An executable image can become input to a binary analysis or translation system.
The loader establishes mapped bytes and entry addresses; a lifter then interprets instruction encodings as explicit machine-state operations.
Symbols and relocation records can help recover intent when retained, but their absence does not authorize guessing missing control-flow targets.
Chapter~\ref{ch:binary-lifting} follows this boundary with a small real x86 subset and a hardware flag oracle.

Its arithmetic model is separate from this linker's implementation.
The linker establishes address and image-format contracts, whereas lifting must preserve instruction-level widths, flags, and partial-register effects.
Cross-architecture translation composes both kinds of obligations with a guest-environment model.

\section{Exercises}
\label{sec:ch60-exercises}
\begin{enumerate}
\item\label{ex:ch60-order} Declare data before code in the Lua region array while leaving addresses unchanged. State the placement order and required program-header order.
\item\label{ex:ch60-size} A region has capacity 4096, used extent 4090, and the next section requires alignment 16 and size 8. Compute the aligned start and explain the rejection.
\item\label{ex:ch60-entry} A symbol is defined at exactly the end of a text section. Why can it be a legitimate boundary symbol but not the entry point accepted by this linker?
\item\label{ex:ch60-reloc-bounds} A section has 12 initialized bytes. Which offsets permit an eight-byte relocation field? State a subtraction-based bounds predicate.
\item\label{ex:ch60-phase} Adding dead-section removal after relocation changes the address of answer. Identify every recorded result in the demo that must be reconsidered before publishing the image.
\end{enumerate}

\section{Further reading}
\label{sec:ch60-reading}
The ELF specification defines the validated graph and runtime image structures~\cite{xinuos-2026-elf}.
The AMD64 supplement fixes the relocation semantics applied to those structures~\cite{matz-2012-amd64}.
Lua's reference manual and GNU linker's script documentation provide contrasting interfaces for programmable layout policy~\cite{ierusalimschy-2026-lua,gnu-2026-ld}.
